% !TeX program = xelatex
% 沉默的代价:条件性披露不足、价格发现与股票截面收益
% 编译: xelatex silence_premium.tex (需 ctex/xeCJK/fandol/booktabs/threeparttable)
\documentclass[UTF8, a4paper, fontset=fandol, zihao=-4]{ctexart}

\usepackage{geometry}
\geometry{left=2.6cm, right=2.6cm, top=2.6cm, bottom=2.6cm}
\usepackage{amsmath, amssymb}
\usepackage{booktabs}
\usepackage{threeparttable}
\usepackage{multirow}
\usepackage{graphicx}
\usepackage{float}
\usepackage{caption}
\usepackage{setspace}
\usepackage{array}
\usepackage[hidelinks]{hyperref}

\captionsetup[table]{labelsep=quad, font=small, skip=4pt}
\captionsetup[figure]{labelsep=quad, font=small, skip=4pt}
\renewcommand{\arraystretch}{1.15}
\onehalfspacing
\setlength{\parindent}{2em}

\ctexset{
  section    = {format = \raggedright\zihao{4}\bfseries, name={,、}, number=\chinese{section}, aftername=},
  subsection = {format = \raggedright\zihao{-4}\bfseries, number=\arabic{section}.\arabic{subsection}},
}

\title{\zihao{3}\textbf{沉默的代价:条件性披露不足、价格发现\\与股票截面收益}}
\author{\zihao{-4} ——来自中国A股ESG信息披露的证据}
\date{}

\begin{document}
\maketitle
\thispagestyle{empty}

\begin{center}
\begin{minipage}{0.92\textwidth}
\zihao{5}
\noindent\textbf{摘\quad 要:}~在ESG信息强制披露准则缺位、评级机构方法分歧显著的制度环境下,企业披露多寡在多大程度上被资本市场定价?本文以2014—2023年A股1{,}339家上市公司为样本,构造\textbf{条件性披露不足}测度:以行业内截面标准化的“ESG披露水平—ESG实际表现”之差刻画企业在同业中的相对披露姿态,并在月度截面上剔除六项风格暴露、行业固定效应及ESG表现与基本面的影响。基于2015年1月至2024年12月的月度截面检验,本文发现:被系统性定价的不是“言过其实”,而是“该说不说”。经完全中性化后的披露裂口对未来一个月收益具有稳定预测力(IC${}=0.0122$,$t=3.45$;防御性口径),披露最不足的十分位组合相对可比公司年化跑输4.17\%($t=-3.14$),经中国三因子调整后为$-4.16$\%($t=-2.92$),且多空组合对市场、规模、价值因子近乎零暴露。

进一步地,本文排除三类竞争性解释并判别两条机制渠道。第一,该效应\textbf{并非漂绿惩罚}:治理维度裂口的预测力与环境维度同样强,高污染行业未见效应放大,处罚公告日的累计异常收益显著为正而非为负;第二,该效应\textbf{并非流动性溢价}:沉默公司反而更具流动性,流动性中介仅解释总效应的7.0\%,追加非流动性中性化后预测力基本不变;第三,\textbf{价格发现速度是更可能的渠道}:沉默折价高度集中于价格延迟较高的子样本,且该差异在低特质波动子样本内依然成立。在组合层面,于全市场等权组合中剔除披露最不足的10\%股票可获得年化0.47\%的超额收益(信息比率1.07),30个基点双边成本后仍保留约九成收益。\textbf{本文结论应理解为对截面可预测性的刻画而非因果推断},且该效应在市值加权口径下消失。

\vspace{4pt}
\noindent\textbf{关键词:}~ESG信息披露;漂绿;价格发现;信息壁垒;截面收益

\noindent\textbf{JEL分类号:}~G12,G14,M41,Q56
\end{minipage}
\end{center}

\vspace{6pt}
\begin{center}
\begin{minipage}{0.92\textwidth}
\zihao{5}
\noindent\textbf{The Cost of Silence: Conditional Under-disclosure, Price Discovery, and the Cross-Section of Stock Returns}

\vspace{4pt}
\noindent\textbf{Abstract:}~Under a voluntary ESG disclosure regime with substantial rating disagreement, which end of the disclosure distribution does the market price? Using 1{,}339 Chinese A-share firms over 2014--2023, we construct a measure of \emph{conditional under-disclosure}: the industry-standardized gap between a firm's ESG disclosure score and its realized ESG performance, orthogonalized each month against six style exposures, industry fixed effects, ESG performance and fundamentals. We find that what is systematically priced is not exaggeration but silence. The orthogonalized disclosure gap predicts one-month-ahead returns (IC${}=0.0122$, $t=3.45$), and the most silent decile underperforms comparable firms by 4.17\% annually ($-4.16$\%, $t=-2.92$, after the Liu-Stambaugh-Yuan China three-factor adjustment), with near-zero factor loadings. We rule out greenwashing punishment (a governance-dimension placebo is equally strong; penalty announcements generate \emph{positive} abnormal returns) and liquidity (silent firms are more liquid; the mediation share is 7.0\%), and provide evidence for a price-discovery channel: the discount concentrates among stocks with high price delay, and survives controlling for idiosyncratic volatility. Excluding the bottom disclosure decile from an equal-weighted portfolio earns 0.47\% per year (IR${}=1.07$), robust to 30bp round-trip costs. Results describe cross-sectional predictability rather than causal effects, and vanish under value weighting.

\vspace{4pt}
\noindent\textbf{Keywords:}~ESG Disclosure; Greenwashing; Price Discovery; Information Barrier; Cross-Section of Returns
\end{minipage}
\end{center}

\newpage
\setcounter{page}{1}

\section{引言}

自2000年联合国全球契约组织提出ESG理念以来,ESG投资已成为全球可持续金融的主流实践。截至2023年,全球签署联合国责任投资原则的机构投资者超过5{,}000家,管理资产规模逾121万亿美元。中国的ESG投资虽起步较晚,但在“双碳”战略推动下发展迅速:截至2025年5月,国内ESG公募基金数量达669只,发行规模连续三年超过2{,}000亿元。2024年4月12日,沪深北三大交易所发布《上市公司可持续发展报告指引》,标志着中国ESG信息披露从完全自愿向部分强制过渡。

然而,在强制统一的可持续披露准则落地之前,中国ESG信息环境长期存在两个结构性问题:一是\textbf{披露的自愿性与主观性},企业对披露内容与详略拥有高度自由裁量权;二是\textbf{评级方法的分歧},国内外评级机构高度依赖企业自愿提供的信息,口径各异。本文数据显示,三家主流机构对A股公司ESG评分的两两秩相关系数仅为0.26至0.69。在这样的信息环境中,一个自然的问题是:\textbf{投资者究竟在为ESG披露的哪一端定价?}

既有文献主要关注披露的“上端”,即漂绿(greenwashing)。Yu等(2020)、Hu等(2023)等研究普遍隐含一个前提:企业夸大其环境绩效的行为会被市场识别并惩罚。这一前提在理论上自然,在中国情境下却缺乏系统的截面定价证据。本文的核心发现与之相反:\textbf{在A股市场上,被系统性定价的不是“言过其实”,而是“该说不说”}。

本文的识别策略建立在一个简单构造之上:企业的绝对披露水平混杂了规模、行业与可见度等大量非信息成分,因此本文首先在行业内将披露水平与实际ESG表现分别标准化并取差,得到“披露裂口”;继而在每个月度截面上将该裂口对六项风格因子、行业哑变量以及ESG表现、盈利能力与杠杆率做正交化,得到\textbf{条件性披露不足}测度。其经济含义是:在规模、风格、行业与真实ESG表现均可比的公司之中,谁在相对地保持沉默。

基于119个月的截面检验,本文得到四组结果。

\textbf{第一,沉默被定价,且不能由经典风险因子解释。} 原始披露裂口不具预测力(IC${}=-0.0004$,$t=-0.10$),但经完全中性化后预测力显著(基准口径IC${}=0.0136$,$t=4.23$;防御性口径0.0122,$t=3.45$)。披露最不足的十分位组合相对可比公司年化跑输4.17\%,经中国三因子调整后为$-4.16$\%($t=-2.92$),多空组合对三因子暴露近乎为零。作为对照,ESG表现水平本身在风格中性化后即失去预测力,说明本文捕捉的是独立于ESG水平的增量信息。

\textbf{第二,这一效应不是漂绿惩罚。} 三重排除检验一致:治理维度裂口的预测力(IC${}=0.0120$,$t=3.62$)不弱于环境维度(0.0099,$t=3.37$),而两者相关系数仅0.18;高污染行业未见一致放大;处罚公告窗内高裂口公司获得显著为正的累计异常收益($+2.76$\%,$t=4.98$)。三项证据共同表明,该信号刻画的是\textbf{跨维度的披露行为},而非绿色维度特有的欺骗行为。

\textbf{第三,机制在于价格发现速度,而非流动性。} 流动性渠道被否定:沉默公司在同规模、同换手条件下反而更具流动性,流动性中介仅解释总效应的7.0\%,追加非流动性中性化后预测力几乎不变。价格发现渠道获得支持:沉默折价高度集中于价格延迟较高的子样本(RankIC${}=0.0168$,$t=3.57$;对照组0.0082,$t=1.66$),且在低特质波动子样本内该判别力依然成立,表明其并非“套利受限”的同义反复。

\textbf{第四,经济意义存在但有限,边界清晰。} 剔除披露最不足10\%股票的等权组合年化超额0.47\%(信息比率1.07),组合增量单边换手仅1.06\%/月,30个基点双边成本后仍保留约九成收益;该信号与质量因子近乎正交(相关系数0.001)。但在市值加权口径下超额消失,其经济价值局限于等权或中小盘倾斜的组合实现。

本文可能的边际贡献有三。\textbf{其一,在测度上},提出“条件性披露不足”的构造思路,并证明正交化处理是信号显现的必要条件,这解释了为何直接使用原始披露或漂绿指标的研究难以得到稳健的定价结论。\textbf{其二,在事实上},以一组相互印证的排除性检验推翻了漂绿文献隐含的“夸大被惩罚”前提。\textbf{其三,在机制上},通过两条渠道的判别检验与壁垒拆除事件研究,将信息壁垒的定价后果落到可观测的价格发现速度上。

\section{制度背景与研究假说}

\subsection{制度背景}

中国ESG信息披露制度的核心特征是\textbf{长期自愿、近期部分强制}。在2024年4月《上市公司可持续发展报告指引》发布之前,除少数行业性规定外,上市公司ESG或社会责任报告的披露与否、详略程度均由企业自主决定。据中国证监会数据,截至2024年底,A股已有2{,}200余家上市公司披露ESG报告或社会责任报告,但披露质量与可比性差异显著。

自愿披露制度带来两个后果:第一,\textbf{披露量本身成为评级的输入},导致“披露多”与“表现好”在评级中难以分离;第二,\textbf{机构间分歧显著}。本文使用的三家机构评分两两秩相关仅为0.26至0.69,与Hu等(2023)刻画的“绿色迷雾”现象一致。这一背景使中国市场成为检验“信息披露如何被定价”的合适场景。

\subsection{理论框架与研究假说}

本文的分析框架建立在信息不完全条件下的资产定价文献之上。Merton(1987)的投资者认知假说指出,当投资者仅了解部分证券时,不为人知的证券需要更高的预期回报以补偿其影子成本;Easley和O'Hara(2004)进一步证明,信息不对称程度较高的股票要求更高的必要回报。两者共同含义是:\textbf{信息壁垒应当被定价}。但定价方向取决于其后果表现为风险补偿还是错误定价。据此提出:

\vspace{2pt}
\noindent\textbf{假说H1(沉默折价):}~在规模、风格、行业与真实ESG表现可比的条件下,相对披露不足的公司未来收益更低。

\vspace{2pt}
\noindent\textbf{假说H2(流动性渠道):}~相对披露不足的公司具有更高的非流动性;流动性差异可解释沉默折价的相当部分。

\vspace{2pt}
\noindent\textbf{假说H3(价格发现渠道):}~相对披露不足的公司价格延迟更高;沉默折价在价格延迟较高的子样本中更强。

\vspace{2pt}
\noindent\textbf{假说H4(壁垒拆除):}~在强制信息释放事件发生时,此前信息壁垒较高的公司应出现显著价格反应。

\vspace{2pt}
需要说明的是,既有漂绿文献隐含的假说是“夸大被惩罚”,即高披露裂口公司未来收益更低,与H1方向相反。本文将其作为\textbf{竞争性假说}予以检验。

\section{研究设计}

\subsection{数据来源与样本}

本文使用三类数据。\textbf{ESG数据}中,披露端为彭博ESG披露评分(衡量披露的广度与详细程度,取值0—100),表现端为国内主流ESG评级的数值化评分;\textbf{财务数据}来自CSMAR数据库,样本区间2014—2023年,经合并与缺失处理后得到9{,}835个公司—年观测,覆盖1{,}339家公司与71个行业。\textbf{行情数据}为A股全市场日频交易数据(2014年1月至2024年12月),用于构造月度收益、风格因子及流动性与价格发现指标。\textbf{评级分歧数据}来自华证、秩鼎与MSCI三家机构。

将年度ESG信号与月度行情合并后,最终样本为2015年1月至2024年12月的120个月度截面(其中119个月具有前瞻收益),共116{,}244个股票—月观测,月均969只股票。

\subsection{条件性披露不足的测度}

第一步,构造行业内披露裂口:
\begin{equation}
GWS_{i,t} = z\!\left(Disclosure_{i,t}\mid Ind_i\right) - z\!\left(Performance_{i,t}\mid Ind_i\right)
\end{equation}
其中$z(\cdot\mid Ind)$表示行业内截面标准化。$GWS$越低,表示企业“做了但不说”的程度越高。

第二步,在每个月度截面上正交化:
\begin{equation}
GWS_{i,t} = \alpha_t + \sum_k \beta_{k,t}\,Style_{k,i,t} + \sum_j \gamma_{j,t}\,Ind_{j,i} + \delta_t'\,Fund_{i,t} + \varepsilon_{i,t}
\end{equation}
风格变量$Style$包括规模(总市值对数)、账面市值比、动量($t-12$至$t-2$月累计收益)、贝塔(过去250个交易日对中证全指回归)、残差波动与换手率,均在截面内标准化并缩尾于$\pm 3$;$Fund$包括ESG表现评分、总资产收益率与资产负债率。残差$\varepsilon_{i,t}$即本文核心变量$GWS^{\perp}$。

\subsection{信息时点处理}

数据核查显示,年度$t$的ESG评分对应会计期间截止日为$t$年12月31日,即基于$t$年年报计算,而A股年报通常于$t+1$年4—6月披露完毕。据此设定两个口径:\textbf{基准口径}将年$t$信号应用于$t+1$自然年的12个月截面;\textbf{防御性口径}(本文主口径)自$t+1$年7月起应用并覆盖其后12个月,确保信号在使用时点必然可得,同时将行业标准化改为行业$\times$年度标准化,消除标准化常数跨年使用的技术性前视。正文主表同时报告两个口径。

\subsection{描述性统计}

表\ref{tab:desc}报告主要变量的描述性统计。披露裂口的标准差为1.21,取值范围$[-4.38, 4.31]$,表明企业间披露姿态存在充分的横截面差异。

\begin{table}[H]
\centering
\begin{threeparttable}
\caption{描述性统计}\label{tab:desc}
\zihao{5}
\begin{tabular}{lrrrrrr}
\toprule
变量 & 观测数 & 均值 & 标准差 & 中位数 & 最小值 & 最大值 \\
\midrule
总资产收益率 & 4{,}290 & 0.0752 & 0.0885 & 0.0693 & $-0.4420$ & 0.2779 \\
披露裂口$GWS$ & 4{,}290 & 0.0000 & 1.2138 & $-0.0056$ & $-4.3766$ & 4.3054 \\
ESG表现评分 & 4{,}290 & 4.5354 & 0.9890 & 4.5000 & 1.0000 & 8.0000 \\
ESG披露评分 & 4{,}290 & 33.9629 & 10.0918 & 31.4706 & 11.1256 & 78.4127 \\
资产负债率 & 4{,}290 & 0.4700 & 0.1711 & 0.4572 & 0.1132 & 0.8215 \\
第一大股东持股 & 4{,}290 & 36.6433 & 14.4752 & 33.0300 & 9.4900 & 72.2300 \\
托宾$Q$ & 4{,}290 & 1.7375 & 1.2951 & 1.3586 & 0.6413 & 29.1670 \\
\bottomrule
\end{tabular}
\begin{tablenotes}\zihao{6}
\item 注:样本为2014—2023年公司—年观测(回归子样本)。
\end{tablenotes}
\end{threeparttable}
\end{table}

\section{沉默折价:基本事实}

\subsection{正交化是信号显现的必要条件}

表\ref{tab:ic}与图\ref{fig:levels}报告不同处理层次下披露裂口对未来一个月收益的预测能力,结果呈现清晰的层次性:\textbf{原始裂口完全不具预测力}(IC${}=-0.0004$,$t=-0.10$);仅做风格中性化后信号出现但不显著(0.0062,$t=1.78$);\textbf{只有在进一步剔除ESG表现水平与基本面后,预测力才充分显现}(0.0136,$t=4.23$)。作为对照,ESG表现评分本身无论是否中性化均不具稳定预测力。这一对照的含义是:ESG水平因子的表面预测力主要来自风格暴露,而披露裂口在剥离风格与ESG水平后仍有独立增量。

\begin{table}[H]
\centering
\begin{threeparttable}
\caption{不同处理层次下的截面预测力}\label{tab:ic}
\zihao{5}
\begin{tabular}{lrrrrrrr}
\toprule
信号 & IC & ICIR & $t$(IC) & RankIC & RankICIR & $t$(RankIC) & 胜率 \\
\midrule
$GWS$ 原始 & $-0.0004$ & $-0.01$ & $-0.10$ & $-0.0052$ & $-0.10$ & $-1.08$ & 48.7\% \\
\quad + 风格中性 & 0.0062 & 0.16 & 1.78 & 0.0037 & 0.10 & 1.06 & 54.6\% \\
\quad + 风格 + ESG/基本面 & \textbf{0.0136} & \textbf{0.39} & \textbf{4.23} & \textbf{0.0124} & \textbf{0.32} & \textbf{3.45} & 64.7\% \\
\midrule
ESG表现评分 原始 & 0.0034 & 0.04 & 0.48 & 0.0096 & 0.12 & 1.26 & 46.2\% \\
\quad + 风格中性 & 0.0046 & 0.11 & 1.21 & 0.0052 & 0.14 & 1.48 & 56.3\% \\
\midrule
防御性口径 & 0.0122 & 0.32 & 3.45 & 0.0103 & 0.27 & 2.87 & 62.8\% \\
\bottomrule
\end{tabular}
\begin{tablenotes}\zihao{6}
\item 注:样本期2015年1月—2024年12月,基准口径$T=119$个月,防御性口径$T=113$个月。IC为Pearson相关系数,RankIC为Spearman秩相关;ICIR为IC均值与其标准差之比;$t$值按$ICIR\times\sqrt{T}$计算。
\end{tablenotes}
\end{threeparttable}
\end{table}

\begin{figure}[H]
\centering
\includegraphics[width=0.86\textwidth]{figs/fig1.png}
\caption{正交化层次与预测力}\label{fig:levels}
\begin{minipage}{0.86\textwidth}\zihao{6}
注:柱高为月度IC均值,括号内为$t$值。灰色柱表示未充分中性化的口径。
\end{minipage}
\end{figure}

\subsection{分组结果}

按$GWS^{\perp}$分为十组(防御性口径),月均前瞻收益呈现明显的\textbf{单侧特征}(表\ref{tab:decile}、图\ref{fig:decile}):最沉默组(G1)月均收益0.264\%,显著低于全样本均值0.612\%,而其余各组差异较小。多空组合(G10$-$G1)年化收益4.44\%,年化波动5.63\%,夏普比率0.79,$t=2.42$,月度胜率63.7\%;沉默组相对可比公司年化$-4.17$\%($t=-3.14$)。

这一形态说明:\textbf{信息主要集中在分布的下尾,而非线性排序},这对组合实现方式具有直接含义(见第八节)。

\begin{table}[H]
\centering
\begin{threeparttable}
\caption{条件性披露不足十分组的月均前瞻收益(防御性口径)}\label{tab:decile}
\zihao{5}
\begin{tabular}{lrrrrrrrrrrr}
\toprule
组别 & G1 & G2 & G3 & G4 & G5 & G6 & G7 & G8 & G9 & G10 & 全样本 \\
\midrule
月均收益(\%) & \textbf{0.264} & 0.407 & 0.606 & 0.637 & 0.651 & 0.661 & 0.627 & 0.788 & 0.847 & 0.634 & 0.612 \\
\bottomrule
\end{tabular}
\begin{tablenotes}\zihao{6}
\item 注:G1为披露最不足组。多空组合(G10$-$G1)年化4.44\%,夏普比率0.79,$t=2.42$,最大回撤$-10.25$\%。
\end{tablenotes}
\end{threeparttable}
\end{table}

\begin{figure}[H]
\centering
\includegraphics[width=0.86\textwidth]{figs/fig2.png}
\caption{十分组月均前瞻收益(防御性口径)}\label{fig:decile}
\end{figure}

\begin{figure}[H]
\centering
\includegraphics[width=0.9\textwidth]{figs/fig3.png}
\caption{累计净值:多空组合与沉默组相对可比公司(防御性口径)}\label{fig:nav}
\end{figure}

\begin{figure}[H]
\centering
\includegraphics[width=0.9\textwidth]{figs/fig4.png}
\caption{月度RankIC序列与12个月滚动均值(防御性口径)}\label{fig:ic}
\end{figure}

\subsection{风险调整后的$\alpha$}

表\ref{tab:alpha}报告组合收益对CAPM、Fama-French三因子及中国三因子(Liu、Stambaugh和Yuan,2019,下称CH-3)的时间序列回归。CH-3因子由本文自行构造:剔除市值最小30\%的股票以规避壳价值污染,价值维度使用盈利价格比,$2\times 3$独立排序后市值加权。样本期内MKT、SMB、VMG的年化收益分别为7.14\%、3.19\%与9.53\%。标准误采用Newey-West调整(滞后6期)。

\begin{table}[H]
\centering
\begin{threeparttable}
\caption{风险调整后的$\alpha$}\label{tab:alpha}
\zihao{5}
\begin{tabular}{llrrrrrr}
\toprule
组合 & 模型 & $\alpha$(年化) & $t$ & $\beta_{MKT}$ & $\beta_{SMB}$ & $\beta_{VMG/HML}$ & $R^2$ \\
\midrule
\multirow{3}{*}{沉默组G1(相对可比公司)}
 & CAPM & $-4.17$\% & $-3.04$ & 0.00 & & & 0.00 \\
 & CH-3 & $\mathbf{-4.16}$\% & $\mathbf{-2.92}$ & 0.02 & $-0.11$ & 0.00 & 0.12 \\
 & FF3 & $-3.42$\% & $-2.65$ & 0.01 & $-0.10$ & $-0.02$ & 0.13 \\
\midrule
\multirow{3}{*}{增强组合(剔除尾部10\%)}
 & CAPM & 0.47\% & 3.03 & 0.00 & & & 0.00 \\
 & CH-3 & \textbf{0.46}\% & \textbf{2.91} & 0.00 & 0.01 & 0.00 & 0.12 \\
 & FF3 & 0.38\% & 2.65 & 0.00 & 0.01 & 0.00 & 0.13 \\
\midrule
\multirow{3}{*}{多空组合(G10$-$G1)}
 & CAPM & 4.49\% & 2.59 & $-0.03$ & & & 0.01 \\
 & CH-3 & 3.88\% & 2.08 & $-0.03$ & 0.15 & 0.05 & 0.10 \\
 & FF3 & 3.27\% & 1.78 & $-0.03$ & 0.13 & 0.06 & 0.12 \\
\midrule
\multicolumn{8}{l}{\emph{基准口径(对照)}} \\
沉默组G1 & CH-3 & $-4.63$\% & $-4.00$ & 0.01 & $-0.17$ & $-0.01$ & 0.28 \\
增强组合 & CH-3 & 0.52\% & 4.00 & 0.00 & 0.02 & 0.00 & 0.28 \\
多空组合 & CH-3 & 5.30\% & 3.38 & $-0.01$ & 0.07 & 0.04 & 0.03 \\
\bottomrule
\end{tabular}
\begin{tablenotes}\zihao{6}
\item 注:主面板为防御性口径($T=113$),下方为基准口径($T=119$)对照。标准误经Newey-West调整(滞后6期)。\textbf{“沉默组”为相对覆盖域等权基准的超额收益};若考察绝对收益,G1的CH-3~$\alpha$仅为$-0.25$\%($t=-0.18$)而G10为3.62\%($t=2.21$),原因是覆盖域(偏大盘)等权组合本身相对CH-3具有正$\alpha$。因此准确表述为:沉默公司相对同类可比公司存在折价,而非其绝对收益为负。
\end{tablenotes}
\end{threeparttable}
\end{table}

\section{竞争性解释的排除}

\subsection{不是漂绿惩罚:E/S/G维度的安慰剂检验}

若市场定价的是“绿色包装”,则效应应集中于环境维度。表\ref{tab:esg}分别以彭博E、S、G分项披露评分构造裂口。结果显示,\textbf{治理维度裂口的预测力不弱于环境维度},而两个分项因子的相关系数仅为0.183——两个近乎正交的信号具有同等预测力,说明捕捉的是跨维度的共同成分。此外,三分项等权合成后的预测力(IC${}=0.0132$,$t=4.18$)与总分裂口基本相同,未见增益。

\begin{table}[H]
\centering
\begin{threeparttable}
\caption{E/S/G分项裂口的预测力(安慰剂检验)}\label{tab:esg}
\zihao{5}
\begin{tabular}{lrrrr}
\toprule
信号 & IC & $t$(IC) & RankIC & $t$(RankIC) \\
\midrule
总分裂口 & 0.0136 & 4.23 & 0.0124 & 3.45 \\
E裂口(环境) & 0.0099 & 3.37 & 0.0067 & 1.90 \\
S裂口(社会) & 0.0066 & 2.27 & 0.0050 & 1.55 \\
G裂口(治理,安慰剂) & \textbf{0.0120} & \textbf{3.62} & \textbf{0.0094} & \textbf{2.75} \\
\bottomrule
\end{tabular}
\begin{tablenotes}\zihao{6}
\item 注:分项因子相关系数中,E与G为0.183,E与S为0.475,S与G为0.120。若效应源于“绿色包装”,治理维度不应具有同等预测力。
\end{tablenotes}
\end{threeparttable}
\end{table}

\subsection{不是漂绿惩罚:高污染行业异质性}

漂绿动机在高污染行业最强。表\ref{tab:poll}使用三种重污染行业划分口径分组检验,未发现一致的效应放大:口径一下高污染组略强,但口径三下高污染组反而不显著。证据不支持“漂绿动机集中于高污染行业”的推论。

\begin{table}[H]
\centering
\begin{threeparttable}
\caption{重污染行业分组检验}\label{tab:poll}
\zihao{5}
\begin{tabular}{llrrrrr}
\toprule
划分口径 & 子样本 & 月均只数 & RankIC & $t$ & 尾部超额(年化) & $t$ \\
\midrule
\multirow{2}{*}{口径一} & 高污染 & 234 & 0.0173 & 2.70 & $-5.70$\% & $-2.24$ \\
 & 非高污染 & 662 & 0.0113 & 2.92 & $-5.20$\% & $-3.22$ \\
\multirow{2}{*}{口径二} & 高污染 & 310 & 0.0134 & 2.52 & $-4.98$\% & $-2.26$ \\
 & 非高污染 & 586 & 0.0122 & 3.05 & $-5.56$\% & $-3.46$ \\
\multirow{2}{*}{口径三} & 高污染 & 143 & 0.0047 & 0.52 & $-1.57$\% & $-0.47$ \\
 & 非高污染 & 753 & 0.0141 & 4.15 & $-5.62$\% & $-4.06$ \\
\bottomrule
\end{tabular}
\end{threeparttable}
\end{table}

\subsection{不是漂绿惩罚:处罚公告的事件研究}

若漂绿在曝光时受到惩罚,则处罚公告日应观察到负向异常收益,且高裂口公司更甚。本文自巨潮资讯网检索2014—2024年标题含“处罚”的公司公告,共得2{,}094条、涉及921家公司(其中环保类73条,证券监管类1{,}561条),按公告前一月的裂口高低分组计算市场调整累计异常收益(表\ref{tab:event}、图\ref{fig:event})。

结果与“漂绿曝光受罚”假说完全相反:高裂口公司在事件窗内获得显著为正的累计异常收益。一个合乎制度事实的解释是,A股处罚公告多与整改完成同时披露(如“关于收到行政处罚决定书及完成整改的公告”),公告本身构成不确定性的解除,信息壁垒被拆除带来的重估效应主导了处罚本身的负面含义。

\begin{table}[H]
\centering
\begin{threeparttable}
\caption{处罚公告事件研究(市场调整累计异常收益)}\label{tab:event}
\zihao{5}
\begin{tabular}{lrrrrr}
\toprule
分组 & 事件数 & CAR$[-5,+20]$ & $t$ & CAR$[0,+5]$ & 负CAR占比 \\
\midrule
高裂口 & 695 & \textbf{$+2.76$\%} & \textbf{4.98} & $+1.04$\% & 46.5\% \\
低裂口 & 655 & $+0.53$\% & 1.04 & $-0.15$\% & 50.8\% \\
\bottomrule
\end{tabular}
\begin{tablenotes}\zihao{6}
\item 注:市场基准为中证全指;分组依据为公告前一月的$GWS^{\perp}$相对当月中位数。
\end{tablenotes}
\end{threeparttable}
\end{table}

\begin{figure}[H]
\centering
\includegraphics[width=0.86\textwidth]{figs/fig6.png}
\caption{处罚公告前后的累计异常收益路径}\label{fig:event}
\begin{minipage}{0.86\textwidth}\zihao{6}
注:阴影为95\%置信区间;基准口径。
\end{minipage}
\end{figure}

\subsection{不是ESG水平因子,也不是质量因子}

表\ref{tab:ic}已显示ESG表现评分在风格中性化后失去预测力。就质量因子而言,本文构造质量合成指标$Q=z(ROA)-z(LEV)-z(IVOL)$(其中$LEV$为资产负债率,$IVOL$为残差波动),其与$GWS^{\perp}$的相关系数仅为\textbf{0.001}。在质量前30\%的组合上叠加披露不足清单(月均剔除29只),边际贡献为年化0.49\%,信息比率0.67($t=2.10$);Fama-MacBeth回归中披露裂口的边际$t$值为3.68,而质量因子自身仅为0.68(表\ref{tab:quality})。

\begin{table}[H]
\centering
\begin{threeparttable}
\caption{质量基底上的边际增量}\label{tab:quality}
\zihao{5}
\begin{tabular}{lrrr}
\toprule
 & 质量基底($Q$前30\%) & 基底+披露清单 & 边际贡献 \\
\midrule
年化超额 & 1.67\% & 2.17\% & \textbf{0.49}\% \\
信息比率 & 0.30 & 0.40 & \textbf{0.67} \\
$t$ & 0.93 & 1.26 & \textbf{2.10} \\
\bottomrule
\end{tabular}
\begin{tablenotes}\zihao{6}
\item 注:$corr(Q, GWS^{\perp})=0.001$。Fama-MacBeth回归中$Q$的$t$值为0.68,$GWS^{\perp}$的边际$t$值为3.68。
\end{tablenotes}
\end{threeparttable}
\end{table}

\subsection{不是规模效应}

双重排序显示,沉默折价在\textbf{大市值半边更强}(RankIC${}=0.0160$,$t=3.63$;小市值半边0.0080,$t=1.57$),尾部年化跑输分别为$-5.11$\%与$-2.96$\%(见表\ref{tab:double})。这排除了“微盘垃圾股”的解释;市值加权组合中超额消失(第八节)源于被剔除标的在市值加权下权重过低的机械原因。

\section{机制检验:流动性还是价格发现?}

\subsection{流动性渠道:被否定}

表\ref{tab:liq}报告H2的检验,三项结果一致否定流动性渠道作为主要机制。\textbf{其一,方向与假说相反}:Fama-MacBeth回归显示披露裂口与Amihud非流动性正相关(系数0.0407,$t=15.69$;控制规模与换手后为0.0414,$t=16.08$),即相对沉默的公司反而更具流动性。\textbf{其二,中介效应微弱}:间接效应仅占总效应的7.0\%(Sobel $z=2.34$,$p=0.019$),且方向对应经典的非流动性溢价而非“沉默$\to$信息风险$\to$更高必要回报”。\textbf{其三,因子存活}:追加$\ln(Amihud)$中性化后预测力基本不变(IC由0.0136降至0.0128,$t=3.90$),再追加价格延迟后仍为0.0128($t=3.91$)。

\begin{table}[H]
\centering
\begin{threeparttable}
\caption{流动性渠道检验}\label{tab:liq}
\zihao{5}
\begin{tabular}{lrr}
\toprule
检验 & 系数/取值 & $t$或$p$ \\
\midrule
$a$路径:$GWS^{\perp}\to \ln(Amihud)$ & 0.0407 & $t=15.69$ \\
$b$路径:$\ln(Amihud)\to$ 收益(控制$GWS^{\perp}$) & 0.00356 & $t=2.36$ \\
$c$:总效应 & 0.00207 & $t=3.87$ \\
$c'$:直接效应 & 0.00194 & $t=3.69$ \\
间接效应$a\times b$(中介占比) & 0.000145\ (\textbf{7.0}\%) & Sobel $z=2.34$,$p=0.019$ \\
\midrule
追加$\ln(Amihud)$中性化后IC & 0.0128 & $t=3.90$ \\
再追加价格延迟后IC & 0.0128 & $t=3.91$ \\
\bottomrule
\end{tabular}
\end{threeparttable}
\end{table}

\subsection{价格发现渠道:获得支持}

本文采用Hou和Moskowitz(2005)的价格延迟测度:以过去52周的周频收益,分别对当期市场收益、以及当期加四阶滞后市场收益回归,取$DELAY=1-R^2_{\text{受限}}/R^2_{\text{非受限}}$,该值越大表示价格对市场信息的反应越慢。样本内该指标均值0.27、中位数0.19,且随市值单调下降(最小市值组0.304,最大市值组0.250),与文献一致。

直接检验的证据偏弱:沉默公司的价格延迟略高(Fama-MacBeth系数$-0.0016$,$t=-1.71$)。但\textbf{条件证据显著}:表\ref{tab:double}与图\ref{fig:double}显示,沉默折价高度集中于价格延迟较高的子样本,尾部年化跑输$-7.39$\%($t=-3.54$)对$-3.01$\%($t=-1.57$)。

\begin{table}[H]
\centering
\begin{threeparttable}
\caption{双重排序:沉默折价在何种环境中最强}\label{tab:double}
\zihao{5}
\begin{tabular}{llrrrr}
\toprule
分组变量 & 子样本 & RankIC & $t$ & 尾部超额(年化) & $t$ \\
\midrule
\multirow{2}{*}{\textbf{价格延迟}} & 高(壁垒高) & \textbf{0.0168} & \textbf{3.57} & $\mathbf{-7.39}$\% & $\mathbf{-3.54}$ \\
 & 低 & 0.0082 & 1.66 & $-3.01$\% & $-1.57$ \\
\multirow{2}{*}{残差波动} & 高 & 0.0164 & 3.25 & $-6.34$\% & $-2.90$ \\
 & 低 & 0.0103 & 2.13 & $-4.66$\% & $-3.20$ \\
\multirow{2}{*}{非流动性} & 高 & 0.0089 & 1.84 & $-6.38$\% & $-3.72$ \\
 & 低 & 0.0136 & 3.24 & $-3.11$\% & $-1.78$ \\
\multirow{2}{*}{市值} & 高 & 0.0160 & 3.63 & $-5.11$\% & $-3.18$ \\
 & 低 & 0.0080 & 1.57 & $-2.96$\% & $-1.60$ \\
\bottomrule
\end{tabular}
\begin{tablenotes}\zihao{6}
\item 注:按各变量的月度中位数二分;尾部超额为最沉默10\%组合相对子样本均值的年化收益。基准口径。
\end{tablenotes}
\end{threeparttable}
\end{table}

\begin{figure}[H]
\centering
\includegraphics[width=0.9\textwidth]{figs/fig5.png}
\caption{沉默折价在不同信息环境下的强度}\label{fig:double}
\end{figure}

\subsection{判别检验:不是套利受限的同义反复}

价格延迟与特质波动存在相关(样本内秩相关0.37),而高特质波动通常被解释为套利受限。为判别两者,本文在\textbf{低残差波动子样本内}重做延迟分组(表\ref{tab:triple})。结果显示价格延迟依然具有判别力,表明\textbf{价格发现速度携带独立于波动率的信息},H3不是套利受限的重新表述。

\begin{table}[H]
\centering
\begin{threeparttable}
\caption{判别检验:低残差波动子样本内的价格延迟分组}\label{tab:triple}
\zihao{5}
\begin{tabular}{lrrrrr}
\toprule
子样本 & 月均只数 & RankIC & $t$ & 尾部超额(年化) & $t$ \\
\midrule
低波动 $\times$ 高延迟 & 233 & \textbf{0.0160} & \textbf{2.37} & $\mathbf{-5.94}$\% & $\mathbf{-3.11}$ \\
低波动 $\times$ 低延迟 & 233 & 0.0085 & 1.34 & $-3.43$\% & $-1.93$ \\
\bottomrule
\end{tabular}
\end{threeparttable}
\end{table}

\subsection{壁垒拆除的价格反应}

第5.3节的事件研究在此获得机制解释:处罚公告作为强制信息释放,构成信息壁垒的瞬时拆除,高裂口公司显著为正的CAR与“不确定性解除$\to$重估”一致。结合H3的结果,本文的机制图景是:\textbf{壁垒抬高$\to$价格发现变慢$\to$折价累积$\to$壁垒拆除$\to$折价修复}。

\section{异质性:评级分歧与指数成分}

Hu等(2023)提出的“绿色迷雾”假说认为,评级机构分歧越大,ESG信息环境越浑浊。本文以三家机构评分的截面百分位排名的跨机构标准差度量分歧度(三家两两秩相关分别为0.685、0.453与0.255,覆盖至少两家评级的公司5{,}142只),按月度中位数分组(表\ref{tab:disag})。

结果呈现有意义的分化:\textbf{排序类信息在低分歧组更强,而尾部折价集中于高分歧组},与绿色迷雾假说一致——在评级机构都难以形成共识之处,最沉默公司的错误定价最深。需强调的是,三家机构评分为当前快照(2026年8月),位于回测窗口之后,本检验依赖“评级分歧是公司层面持久属性”的假设,结论应作为\textbf{暗示性证据}。

按指数成分分组的结果亦具参考价值:沉默折价的尾部效应在中证500成分内显著($-4.83$\%,$t=-2.89$),在沪深300成分内不显著($-2.56$\%,$t=-1.14$),在指数外样本中最强($-6.83$\%,$t=-3.14$)。

\begin{table}[H]
\centering
\begin{threeparttable}
\caption{评级分歧度与指数成分的调节作用}\label{tab:disag}
\zihao{5}
\begin{tabular}{llrrrrr}
\toprule
分组维度 & 子样本 & 月均只数 & RankIC & $t$ & 尾部超额(年化) & $t$ \\
\midrule
\multirow{2}{*}{评级分歧度} & 低分歧 & 445 & \textbf{0.0177} & \textbf{4.35} & $-3.81$\% & $-1.76$ \\
 & 高分歧 & 444 & 0.0088 & 1.72 & $\mathbf{-6.34}$\% & $\mathbf{-3.93}$ \\
\midrule
\multirow{3}{*}{指数成分} & 沪深300内 & 191 & 0.0137 & 2.26 & $-2.56$\% & $-1.14$ \\
 & 中证500内 & 321 & 0.0137 & 2.44 & $-4.83$\% & $-2.89$ \\
 & 指数外 & 384 & 0.0096 & 1.69 & $\mathbf{-6.83}$\% & $\mathbf{-3.14}$ \\
\bottomrule
\end{tabular}
\end{threeparttable}
\end{table}

\section{经济意义}

鉴于信息集中于分布下尾,本文采用\textbf{旗标式}而非线性打分的组合实现:在覆盖域等权组合中剔除披露最不足的10\%股票(月均剔除89只),结果见表\ref{tab:port}。

三点值得强调。\textbf{第一,成本几乎不构成约束}:由于清单仅覆盖组合的十分之一,组合层面的增量单边换手仅1.06\%/月,按30个基点双边成本计费后仍保留约九成超额收益。\textbf{第二,年度冻结版本可行}:一年确定一次名单、清单月换手仅6.3\%的实现方式仍保留0.41\%的年化超额($t=2.43$)。\textbf{第三,存在明确的容量边界}:市值加权口径下超额消失($-0.19$\%,$t=-0.47$),因此该信号的经济价值局限于等权或中小盘倾斜的组合实现,不适用于大资金的市值加权组合。

\begin{table}[H]
\centering
\begin{threeparttable}
\caption{组合实现结果(等权,月度调仓,基准口径)}\label{tab:port}
\zihao{5}
\begin{tabular}{lrrrr}
\toprule
指标 & 月度刷新清单 & 年度冻结清单 & 30bp成本后 & 市值加权 \\
\midrule
年化超额 & 0.57\% & 0.41\% & 0.53\% & $-0.19$\% \\
跟踪误差 & 0.49\% & 0.53\% & 0.49\% & 1.30\% \\
信息比率 & 1.17 & 0.77 & 1.09 & $-0.15$ \\
$t$ & 3.67 & 2.43 & 3.44 & $-0.47$ \\
月度胜率 & 67.2\% & 63.0\% & 66.4\% & 51.3\% \\
盈亏比 & 1.19 & 1.07 & — & — \\
盈利因子 & 2.43 & 1.82 & — & — \\
最大相对回撤 & $-0.44$\% & $-0.68$\% & — & — \\
清单月均换手 & 17.9\% & 6.3\% & — & — \\
组合增量单边换手 & 1.06\%/月 & 未单独测算 & 1.06\%/月 & — \\
\bottomrule
\end{tabular}
\begin{tablenotes}\zihao{6}
\item 注:防御性口径下月度刷新清单的年化超额为0.47\%,信息比率1.07($t=3.28$)。
\end{tablenotes}
\end{threeparttable}
\end{table}

\section{稳健性检验}

\textbf{信号时点。} 将信号应用起点由$t+1$年1月推迟至7月,RankIC由0.0124($t=3.45$)降至0.0103($t=2.87$),增强组合超额由0.57\%降至0.47\%,衰减约15\%,结论不变。

\textbf{标准化口径。} 将行业标准化改为行业$\times$年度标准化,消除跨年常数,结论不变(已含于防御性口径)。

\textbf{信号新鲜度。} 按信号使用月龄分组,前6个月RankIC为0.0123($t=2.30$),后6个月为0.0125($t=2.58$),未见衰减,表明该信号刻画的是缓慢变化的公司属性而非事件型信息。

\textbf{替代构造。} 直接以披露评分残差(不经裂口构造)构造的信号与本文核心变量相关系数达0.971,预测力几乎相同(IC${}=0.0142$,$t=4.21$)。这一结果具有方法论含义:经完全中性化后,“披露—表现裂口”实质上等价于\textbf{条件性披露水平},裂口构造中的表现端主要起到正交化作用。

\textbf{实体绩效证据。} 在年度公司面板中(行业与年度双向固定效应,聚类稳健标准误),披露裂口对总资产收益率的系数在控制ESG表现后为0.0111($t=3.73$),前瞻一期与1\%/99\%缩尾后结论不变;而不控制ESG表现时系数为$-0.0042$($t=-2.23$),呈现典型的遮掩效应。

\textbf{未能实施的识别性检验。} 本文曾尝试以2024年4月《上市公司可持续发展报告指引》作为信息壁垒外生拆除的准自然实验进行双重差分检验,但因三方面原因判定该检验在当前数据下不可识别并予以放弃:(1)新“国九条”于\textbf{同日}(2024年4月12日)发布并引发市场风格剧烈分化,而处理组按市值特征构造,与该冲击的受益方高度重合,两者无法分离;(2)双重差分窗口恰为本文信号的失效期,基线折价在窗口内已反转为正,缺乏可供检验“收窄”的基准效应;(3)三重差分系数不显著(年化2.2\%,$t=0.21$),分市值层估计在$-49$\%至$+38$\%间剧烈波动。相关代码与结果作为否定性记录保留。

\section{结论、局限与展望}

\subsection{主要结论}

本文以2014—2023年A股上市公司为样本,考察ESG信息披露在截面收益中的定价。核心结论是:\textbf{在中国市场,被系统性定价的不是“言过其实”,而是“该说不说”}。在规模、风格、行业与真实ESG表现均可比的条件下,相对披露不足的公司未来收益显著更低,该折价年化约4.2\%(CH-3调整,$t=-2.92$),不能由市场、规模、价值因子解释,亦非流动性溢价、质量因子或规模效应的伪装。机制检验表明,该折价的渠道并非流动性,而是\textbf{价格发现速度}:折价高度集中于价格延迟较高的股票,且该集中性独立于特质波动;以监管处罚公告作为壁垒拆除时点,此前披露裂口较高的公司获得显著正向的累计异常收益。

在方法上,本文强调\textbf{正交化处理是信号显现的必要条件}——未经处理的原始披露裂口不具任何预测力。这解释了以往直接使用披露或漂绿指标的研究难以获得稳健定价结论的原因。

\subsection{局限}

\textbf{第一,本文是可预测性研究,而非因果推断。} 企业披露决策是内生选择,本文的中性化处理仅控制可观测特征,无法排除不可观测的公司特质同时驱动披露与收益。文中一律使用“预测”“定价”而非“导致”的表述,第九节说明的准自然实验尝试未能通过识别性检验。

\textbf{第二,经济量级有限。} 月度IC约0.012、组合年化超额约0.5\%。本文强调的是统计稳健性与机制清晰度,而非收益幅度;该信号更适合作为多因子体系中的低相关增量成分或负面清单,而非独立策略。

\textbf{第三,样本覆盖受限于评级数据。} 披露评分月均覆盖969只股票,偏向大中盘,约占A股全市场五分之一,结论向未覆盖公司的外推需谨慎。

\textbf{第四,测度依赖单一评级源。} 披露端与表现端各来自单一机构,评级方法本身的偏误可能传导至测度;评级分歧度检验使用的是当前快照而非时点数据。

\textbf{第五,近期存在衰减迹象。} 12个月滚动RankIC均值在2023年下半年后转负,2024年多空组合年化收益为$-1.7$\%,可能与强制披露监管趋严导致披露裂口收敛有关。以24个月滚动窗口衡量的衰减尚未触发,现有样本不足以判定该效应是否已系统性消失。

\subsection{展望}

三个方向值得后续研究。\textbf{其一,扩展覆盖}:采用全覆盖的本土评级数据重做检验,特别是考察完全不披露公司(扩展边际)的定价。\textbf{其二,真实的识别策略}:待2025年度报告强制披露落地后,以实际披露行为的变化而非成分身份代理重做双重差分,届时亦无同期政策冲击的污染。\textbf{其三,披露内容的文本刻画}:本文测度基于披露评分,未触及披露文本的具体性与承诺兑现率等更精细维度。

\vspace{10pt}
\section*{参考文献}
\zihao{5}
\begin{enumerate}\setlength{\itemsep}{2pt}
\item Amihud Y. Illiquidity and stock returns: cross-section and time-series effects[J]. Journal of Financial Markets, 2002, 5(1): 31--56.
\item Easley D, O'Hara M. Information and the cost of capital[J]. Journal of Finance, 2004, 59(4): 1553--1583.
\item Fama E F, MacBeth J D. Risk, return, and equilibrium: empirical tests[J]. Journal of Political Economy, 1973, 81(3): 607--636.
\item Hou K, Moskowitz T J. Market frictions, price delay, and the cross-section of expected returns[J]. Review of Financial Studies, 2005, 18(3): 981--1020.
\item Hu X, Hua R, Liu Q, et al. The green fog: environmental rating disagreement and corporate greenwashing[J]. Pacific-Basin Finance Journal, 2023, 78: 101952.
\item Liu J, Stambaugh R F, Yuan Y. Size and value in China[J]. Journal of Financial Economics, 2019, 134(1): 48--69.
\item Merton R C. A simple model of capital market equilibrium with incomplete information[J]. Journal of Finance, 1987, 42(3): 483--510.
\item Yu E P Y, Van Luu B, Chen C H. Greenwashing in environmental, social and governance disclosures[J]. Research in International Business and Finance, 2020, 52: 101192.
\end{enumerate}

\end{document}
